  \documentclass[11pt]{article}
  \usepackage[margin=1in]{geometry}
  \usepackage{amsmath,amssymb,booktabs,array,longtable,enumitem}
  \usepackage[hidelinks]{hyperref}
  
  \newcommand{\E}{\mathbb{E}}
  \newcommand{\Var}{\operatorname{Var}}
  \newcommand{\Cov}{\operatorname{Cov}}
  \newcommand{\R}{\mathbb{R}}
  \newcommand{\ind}{\mathbf{1}}
  \newcommand{\thetav}{\boldsymbol{\theta}}
  
  \title{An AI Productivity Atlas, Prediction-Powered Inference,\\
  and Multi-Task Prediction-Powered Inference}
  \author{}
  \date{\today}
  
  \begin{document}
  \maketitle
  
  \section*{Purpose and source convention}
  
  This is a new, self-contained note. It does not replace or modify the earlier
  AI-project discussion in \texttt{todo\_chat\_merged\_notes.tex}.
  
  The three source documents are:
  \begin{enumerate}[leftmargin=2em]
      \item \emph{The Distributional Effects of AI's Productivity Gains: A Data
      Fusion Approach}, the five-page project proposal (``Proposal'');
      \item Angelopoulos, Bates, Fannjiang, Jordan, and Zrnic,
      \emph{Prediction-Powered Inference}, arXiv:2301.09633v4 (``PPI'');
      \item Emmenegger, Stahler, and Podimata,
      \emph{Prediction-Powered Inference Across Many Tasks for AI Evaluation
      \& Social Science Research}, arXiv:2605.29249v1 (``MT-PPI'').
  \end{enumerate}
  Page citations refer to PDF pages. A statement marked ``derived'' is proved in
  this note from the displayed definitions. A statement marked ``proposal
  claim'' is asserted in the proposal but is not proved there.
  
  \section{The project proposal}
  
  \subsection*{Notation}
  
  \begin{longtable}{>{\raggedright\arraybackslash}p{0.22\linewidth}
                    >{\raggedright\arraybackslash}p{0.70\linewidth}}
  \toprule
  Symbol & Meaning \\
  \midrule
  \endfirsthead
  \toprule
  Symbol & Meaning \\
  \midrule
  \endhead
  \(j\in\{1,\ldots,J\}\) & Occupation--task cell and the number of cells. \\
  \(\theta_j\) & True productivity gain in cell \(j\). \\
  \(\thetav\) & Full map \((\theta_1,\ldots,\theta_J)'\). \\
  \(p_j\) & Economic weight for cell \(j\), with
  \(p_j\geq 0\) and \(\sum_j p_j=1\). \\
  \(\bar\theta\) & Weighted mean gain, \(\sum_jp_j\theta_j\). \\
  \(p_j^0,p_j^A\) & Weights for the pre-AI basket and the basket of tasks
  actually brought to AI. \\
  \(\bar\theta^0,\bar\theta^A\) & Mean gains under the two baskets,
  \(\sum_jp_j^0\theta_j\) and \(\sum_jp_j^A\theta_j\). \\
  \(k\) & Number of cells selected in a top-\(k\) ranking or targeting set. \\
  \(S_j\) & Broad but biased and noisy model score for cell \(j\), after fixing
  one basket. \\
  \(\mu,\tau^2\) & Mean and variance of the Gaussian distribution of
  true cell gains. \\
  \(a,b\) & Intercept and slope of the score-to-truth measurement channel;
  \(0<b<1\) represents compression. \\
  \(\eta_j,\sigma_S^2\) & Score error and its variance. \\
  \(R_j\) & Indicator that cell \(j\) is selected for validation. \\
  \(m\) & Number of validated cells, \(m=\sum_jR_j\). \\
  \(\widehat v_j\) & Experimental estimate of \(\theta_j\) in a validated
  cell. \\
  \(\varepsilon_j,s_j^2\) & Experimental estimation error and its known design
  variance. \\
  \(m_j,\omega^2\) & Posterior mean and variance of \(\theta_j\) given only
  the broad score. \\
  \(P_j,m_j^{\mathrm{val}}\) & Posterior precision and mean when cell \(j\)
  is also validated. \\
  \(\lambda\) & Score reliability ratio,
  \(b^2\tau^2/(b^2\tau^2+\sigma_S^2)\). \\
  \(\psi\) & Score-channel parameter vector
  \((\mu,\tau^2,a,b,\sigma_S^2)\). \\
  \(\Lambda_j\) & Sensitivity constant in the proposal's asserted
  near-oracle excess-risk expansion. \\
  \(\sigma_\theta^2,\beta,\nu_j,\widehat\theta_j^v\) & Source-PDF aliases for
  \(\tau^2,b,\eta_j,\widehat v_j\); they are listed only to make the notation
  conversion explicit. \\
  \bottomrule
  \end{longtable}
  
  \noindent\textit{Notation conversion.}
  The proposal PDF writes the latent variance, channel slope, score error, and
  validation estimate as \(\sigma_\theta^2,\beta,\nu_j,\widehat\theta_j^v\).
  The earlier AI-project note writes the same objects as
  \(\tau^2,b,\eta_j,\widehat v_j\), with validation error
  \(\varepsilon_j\). This note preserves the earlier note's notation.
  
  \subsection*{Board 1: the objects of interest}
  
  Fix one task basket. The primitive object is the vector
  \[
      \thetav
      =(\theta_1,\ldots,\theta_J)'.
  \]
  The mean is
  \[
      \bar\theta=\sum_{j=1}^Jp_j\theta_j.
  \]
  The corresponding weighted cross-cell dispersion is
  \[
      \sum_{j=1}^Jp_j(\theta_j-\bar\theta)^2.
  \]
  The two basket means use the same notation as the earlier note:
  \[
      \bar\theta^0=\sum_{j=1}^Jp_j^0\theta_j,
      \qquad
      \bar\theta^A=\sum_{j=1}^Jp_j^A\theta_j.
  \]
  Rankings and top-\(k\) sets are other functions of the same vector.
  These are the proposal's estimands. An atlas estimator is a rule for estimating
  \(\thetav\). The estimand and the estimator are not the same
  object. (Proposal, p.~1.)
  
  Recovering the full map is sufficient for computing every summary. It is not
  necessary for every summary. In particular, a randomized sample of validated
  cells can identify a weighted mean directly, without identifying every
  \(\theta_j\). Section 4 derives that estimator. Thus, the proposal's statement
  that credible recovery of every summary requires the full map is too strong.
  The full map is essential for cell-level claims and rankings; it is not
  essential for the mean.
  
  \subsection*{Board 2: why scores alone do not identify the map}
  
  The proposal assumes
  \[
      \theta_j\sim N(\mu,\tau^2),
      \qquad
      S_j=a+b\theta_j+\eta_j,
      \qquad
      \eta_j\sim N(0,\sigma_S^2),
  \]
  with \(\eta_j\) independent of \(\theta_j\). (Proposal, p.~4.)
  
  The score distribution reveals only
  \[
      \E[S_j]=a+b\mu
  \]
  and
  \[
      \Var(S_j)=b^2\tau^2+\sigma_S^2.
  \]
  There are five unknown channel parameters
  \[
      (\mu,\tau^2,a,b,\sigma_S^2)
  \]
  which the earlier note collects as
  \(\psi=(\mu,\tau^2,a,b,\sigma_S^2)\),
  but only two marginal score moments. Therefore the score distribution alone
  does not identify the level \(\mu\), the scale \(\tau^2\), or the
  channel parameters separately. This is the proposal's nonidentification
  argument. (Proposal, p.~4; algebra derived.)
  
  The presence of \(\eta_j\) also means that the ranking of the observed scores
  need not equal the ranking of the true gains. The scores identify their own
  ranking. They do not identify the realized ranking of
  \((\theta_1,\ldots,\theta_J)\) unless additional assumptions remove or
  discipline score noise. This qualification is not explicit in the proposal.
  
  \subsection*{Board 3: what randomized validation identifies}
  
  For a validated cell, the proposal observes
  \[
      \widehat v_j=\theta_j+\varepsilon_j,
      \qquad
      \E[\varepsilon_j\mid\theta_j,S_j]=0,
      \qquad
      \Var(\varepsilon_j\mid\theta_j,S_j)=s_j^2.
  \]
  Assume the validation cells are a random sample of the target cell population,
  the design variances \(s_j^2\) are known, and \(\varepsilon_j\), \(\eta_j\), and
  \(\theta_j\) obey the independence conditions in the Gaussian channel. These
  conditions make the validation moments equal to the population moments.
  (Proposal, pp.~2 and 4.)
  
  The five channel parameters are then recovered in five linear steps:
  \begin{align*}
      \mu
      &=\E[\widehat v_j],\\
      \tau^2
      &=\Var(\widehat v_j)-\E[s_j^2],\\
      b
      &=\frac{\Cov(S_j,\widehat v_j)}
              {\tau^2},\\
      a
      &=\E[S_j]-b\mu,\\
      \sigma_S^2
      &=\Var(S_j)-b^2\tau^2.
  \end{align*}
  The third line requires \(\tau^2>0\). Every displayed equality is
  derived from the two measurement equations. This proves the identification
  claim under the stated sampling and independence conditions. The proposal
  states the conclusion but does not show these steps. (Proposal, p.~4.)
  
  \subsection*{Board 4: the empirical-Bayes map}
  
  For an unvalidated cell, Gaussian conditioning gives
  \[
      m_j
      =\E[\theta_j\mid S_j]
      =\mu+
        \frac{b\tau^2}
             {b^2\tau^2+\sigma_S^2}
        (S_j-a-b\mu).
  \]
  The posterior variance and the score reliability ratio are
  \[
      \omega^2
      =
      \frac{\tau^2\sigma_S^2}{b^2\tau^2+\sigma_S^2},
      \qquad
      \lambda
      =
      \frac{b^2\tau^2}{b^2\tau^2+\sigma_S^2}.
  \]
  The coefficient multiplying the centered score is the calibrated shrinkage
  coefficient. It corrects both compression and score noise. This formula is
  derived from the proposal's Gaussian channel.
  
  For a validated cell, both measurements enter. Gaussian precision addition
  gives
  \begin{align*}
      P_j
      &=
        \frac{1}{\tau^2}
        +\frac{b^2}{\sigma_S^2}
        +\frac{1}{s_j^2}
      ,\\
      m_j^{\mathrm{val}}
      &=
      P_j^{-1}\left(
        \frac{\mu}{\tau^2}
        +\frac{b(S_j-a)}{\sigma_S^2}
        +\frac{\widehat v_j}{s_j^2}
      \right).
  \end{align*}
  For an unvalidated cell, remove the last precision and numerator terms.
  These formulas make the proposal's phrase ``posterior means combine the two''
  explicit. (Proposal, pp.~2--4; formulas derived.)
  
  \subsection*{Board 5: what is established and what remains a research claim}
  
  \begin{enumerate}[leftmargin=2em]
      \item Score-only nonidentification follows directly from the displayed
      moments.
      \item Channel identification follows from randomized validation plus the
      measurement-error conditions displayed above.
      \item Cell-level recovery outside the validation sample is model-based:
      it uses the common Gaussian channel and exchangeable distribution of
      \(\theta_j\).
      \item The statement that feasible map risk exceeds oracle risk by
      \(\Lambda_j/m+o(m^{-1})\), with an aggregate error of order
      \(m^{-1/2}\), is a proposal claim. The five-page proposal supplies no
      proof or regularity conditions. (Proposal, p.~4.)
      \item The numerical ``knee'' near \(m=40\) is a simulation result, not an
      identification theorem. (Proposal, pp.~2--3.)
      \item The claim that true value gain lies between the pre-AI and actual-use
      basket gains needs an ordering or behavioral assumption connecting the two
      baskets. The proposal states the bound but does not supply that assumption
      or its proof. (Proposal, pp.~1 and 3.)
  \end{enumerate}
  
  \section{Angelopoulos et al.: prediction-powered inference}
  
  \subsection*{Notation}
  
  \begin{longtable}{>{\raggedright\arraybackslash}p{0.22\linewidth}
                    >{\raggedright\arraybackslash}p{0.70\linewidth}}
  \toprule
  Symbol & Meaning \\
  \midrule
  \endfirsthead
  \toprule
  Symbol & Meaning \\
  \midrule
  \endhead
  \(i\) & Observation index. \\
  \(n,N\) & Numbers of labeled and unlabeled observations in the
  superpopulation setup. \\
  \((X_i,Y_i)\) & Labeled feature--outcome pair. \\
  \((\widetilde X_i,\widetilde Y_i)\) & Unlabeled feature--outcome pair;
  \(\widetilde Y_i\) is not observed. \\
  \(f\) & Prediction rule; \(f(X)\) predicts \(Y\). \\
  \(\theta^\star\) & Estimand. For mean estimation,
  \(\theta^\star=\E[Y_i]\). \\
  \(\theta\) & Generic candidate value used in the loss, subgradient, and
  confidence-set construction. \\
  \(\widehat\theta_{\mathrm{class}}\) & Labeled-only sample mean. \\
  \(\widetilde\theta_f\) & Imputed estimate computed from predictions. \\
  \(\Delta,\widehat\Delta\) & Population rectifier and its labeled-sample
  estimate in the mean problem. \\
  \(\widehat\theta_{\mathrm{PP}}\) & Prediction-powered point estimator of the
  mean. \\
  \(\sigma_Y^2,\sigma_{f-Y}^2,\sigma_f^2\) & Population variances of \(Y\),
  \(f(X)-Y\), and \(f(X)\); hats denote their sample estimates. \\
  \(\alpha,\delta\) & Total error probability and the part assigned to the
  rectifier confidence set. \\
  \(z_{1-\alpha/2}\) & Standard-normal quantile used in a two-sided interval. \\
  \(\ell_\theta\) & Convex loss indexed by \(\theta\). \\
  \(p\) & Dimension of the parameter \(\theta\in\R^p\). \\
  \(g_\theta\) & Subgradient of \(\ell_\theta\) with respect to \(\theta\). \\
  \(\Delta_\theta\) & General rectifier
  \(\E[g_\theta(X,Y)-g_\theta(X,f(X))]\). \\
  \(g_\theta^f\) & Imputed population subgradient
  \(\E[g_\theta(X,f(X))]\). \\
  \(R_\delta(\theta)\) & Confidence set for \(\Delta_\theta\). \\
  \(T_{\alpha-\delta}(\theta)\) & Confidence set for \(g_\theta^f\). \\
  \(C_{\mathrm{PP}}\) & Prediction-powered confidence set for
  \(\theta^\star\). \\
  \(L\) & Uniformly sampled labeled subset of a fixed population. \\
  \(\theta_N\) & Fixed-population mean, \(N^{-1}\sum_iY_i\). \\
  \(\widehat\theta_{\mathrm{PP},N}\) & Finite-population PPI estimator. \\
  \(e_i,S_e^2\) & Finite-population residual \(Y_i-f(X_i)\) and its finite-
  population sample variance. \\
  \bottomrule
  \end{longtable}
  
  \subsection*{Board 1: estimand, estimator, and inferential object}
  
  For mean estimation, the estimand is
  \[
      \theta^\star=\E[Y_i].
  \]
  The estimator is a random function of observed data:
  \[
      \widehat\theta_{\mathrm{PP}}
      =
      \frac{1}{N}\sum_{i=1}^N f(\widetilde X_i)
      -
      \frac{1}{n}\sum_{i=1}^n\{f(X_i)-Y_i\}.
  \]
  The inferential object is a confidence interval centered on that estimator.
  The paper is interested in the scientific estimand \(\theta^\star\); PPI is the
  method used to estimate it and quantify uncertainty. (PPI, pp.~2--3.)
  
  \subsection*{Board 2: why the correction works}
  
  Rewrite the estimator as
  \[
      \widehat\theta_{\mathrm{PP}}
      =
      \underbrace{\frac{1}{N}\sum_{i=1}^N f(\widetilde X_i)}
                  _{\text{many cheap predictions}}
      +
      \underbrace{\frac{1}{n}\sum_{i=1}^n\{Y_i-f(X_i)\}}
                  _{\text{mean prediction error on labels}}.
  \]
  Under the paper's initial setup, the labeled and unlabeled samples are
  independent draws from the same population and \(f\) is independent of those
  samples. Therefore
  \begin{align*}
      \E[\widehat\theta_{\mathrm{PP}}]
      &=
      \E[f(X_i)]+\E[Y_i-f(X_i)]\\
      &=
      \E[Y_i]\\
      &=\theta^\star.
  \end{align*}
  No line requires \(f\) to be unbiased. Systematic prediction bias appears in
  both prediction terms and cancels. (PPI, pp.~2--3; algebra derived.)
  
  Because the two samples are independent,
  \[
      \Var(\widehat\theta_{\mathrm{PP}})
      =
      \frac{\sigma_f^2}{N}
      +
      \frac{\sigma_{f-Y}^2}{n}.
  \]
  The corresponding large-sample interval is
  \[
      \widehat\theta_{\mathrm{PP}}
      \ \pm\
      z_{1-\alpha/2}
      \sqrt{\frac{\widehat\sigma_f^2}{N}
            +\frac{\widehat\sigma_{f-Y}^2}{n}}.
  \]
  The labeled-only estimator has variance \(\sigma_Y^2/n\). When \(N\) is much
  larger than \(n\), PPI is more precise when the prediction residual has lower
  variance than the outcome:
  \[
      \sigma_{f-Y}^2<\sigma_Y^2.
  \]
  Accurate predictions improve precision; correction, not predictive accuracy,
  provides validity. (PPI, pp.~3--4 and 36.)
  
  \subsection*{Board 3: the general rectifier}
  
  The paper considers estimands defined by
  \[
      \theta^\star
      =\arg\min_{\theta\in\R^p}\E[\ell_\theta(X_i,Y_i)].
  \]
  For a nondegenerate convex problem, the minimizer satisfies
  \[
      \E[g_{\theta^\star}(X_i,Y_i)]=0.
  \]
  Add and subtract the imputed subgradient:
  \begin{align*}
      \E[g_\theta(X_i,Y_i)]
      &=
      \E[g_\theta(X_i,f(X_i))]\\
      &\quad+
      \E[g_\theta(X_i,Y_i)-g_\theta(X_i,f(X_i))]\\
      &=g_\theta^f+\Delta_\theta.
  \end{align*}
  The first term is estimated precisely with many predictions. The second term is
  estimated with the labels. If
  \[
      \Pr\{\Delta_\theta\in R_\delta(\theta)\}\geq1-\delta
  \]
  and
  \[
      \Pr\{g_\theta^f\in T_{\alpha-\delta}(\theta)\}
      \geq1-(\alpha-\delta),
  \]
  then the paper defines
  \[
      C_{\mathrm{PP}}
      =
      \{\theta:0\in R_\delta(\theta)+T_{\alpha-\delta}(\theta)\}.
  \]
  By the union bound, \(\theta^\star\) belongs to this set with probability at
  least \(1-\alpha\). This is the paper's main convex-estimation theorem.
  (PPI, pp.~4--5.)
  
  \subsection*{Board 4: the fixed-population version}
  
  Suppose the \(N\) units are fixed, \(f(X_i)\) is observed for every unit, and
  \(L\) is a simple random sample without replacement of \(n\) units
  whose \(Y_i\)'s are observed. The finite-population estimand is
  \[
      \theta_N=\frac{1}{N}\sum_{i=1}^NY_i.
  \]
  The natural finite-population PPI estimator is
  \[
      \widehat\theta_{\mathrm{PP},N}
      =
      \frac{1}{N}\sum_{i=1}^Nf(X_i)
      +
      \frac{1}{n}\sum_{i\in L}\{Y_i-f(X_i)\}.
  \]
  Let \(e_i=Y_i-f(X_i)\). Then
  \begin{align*}
      \widehat\theta_{\mathrm{PP},N}-\theta_N
      &=
      \frac{1}{n}\sum_{i\in L}e_i
      -
      \frac{1}{N}\sum_{i=1}^Ne_i.
  \end{align*}
  Simple random sampling makes the sample residual mean unbiased for the
  population residual mean. Hence the estimator is exactly design-unbiased.
  Its design variance is
  \[
      \Var(\widehat\theta_{\mathrm{PP},N})
      =
      \frac{1}{n}\left(1-\frac{n}{N}\right)S_e^2.
  \]
  This is the fixed-population difference estimator. The paper states the
  finite-population rectifier construction; the displayed point estimator and
  variance follow by applying its sampling-without-replacement formula to the
  residuals. (PPI, pp.~23--24 and 34; last two displays derived.)
  
  \subsection*{Board 5: what PPI does and does not do}
  
  \begin{enumerate}[leftmargin=2em]
      \item PPI gives valid inference for a specified low-dimensional estimand,
      such as a mean, quantile, or regression coefficient. It does not require a
      correct prediction model. (PPI, pp.~1 and 5--8.)
      \item Validity requires the labeled observations to represent the target
      population under the sampling assumptions. Prediction correction cannot
      repair a nonrepresentative label design without additional adjustment.
      (PPI, pp.~3, 15--16, and 23.)
      \item A single aggregate can be identified without estimating every
      latent unit outcome.
      \item PPI does not, by itself, recover every unobserved cell effect
      \(\theta_j\). It corrects a functional whose average prediction error is
      learnable from the labels.
      \item In the fixed-population case, the only sampling uncertainty comes
      from estimating the population-average residual. The population mean of
      the predictions is already observed. (PPI, p.~23; derived above.)
  \end{enumerate}
  
  \section{Emmenegger, Stahler, and Podimata: multi-task PPI}
  
  \subsection*{Notation}
  
  \begin{longtable}{>{\raggedright\arraybackslash}p{0.22\linewidth}
                    >{\raggedright\arraybackslash}p{0.70\linewidth}}
  \toprule
  Symbol & Meaning \\
  \midrule
  \endfirsthead
  \toprule
  Symbol & Meaning \\
  \midrule
  \endhead
  \(\mathcal T,t,T\) & Set of tasks, a task index, and number of tasks. \\
  \(i,N\) & Unit index and fixed number of units per task. \\
  \(X_i^{(t)},Y_i^{(t)}\) & Covariates and ground-truth outcome for unit \(i\)
  in task \(t\). \\
  \(\widehat Y_i^{(t)}\) & Proxy outcome observed for every unit in task \(t\). \\
  \(\theta^{(t)}\) & Task-\(t\) finite-population mean. \\
  \(L^{(t)},n_t,O_i^{(t)}\) & Simple-random labeled subset, its size, and the
  indicator \(\ind\{i\in L^{(t)}\}\). \\
  \(s\) & Recalibration map applied to the proxy; \(s_i\) abbreviates
  \(s(X_i,\widehat Y_i)\) in a single task. \\
  \(\overline Y_L,\overline Y_N\) & Labeled-sample and population means of
  \(Y\) in a single task. \\
  \(\overline s_L,\overline s_N\) & Labeled-sample and population means of the
  recalibrated proxy. \\
  \(\gamma_t,\gamma^\star(s),\widehat\gamma_t\) & Task-specific power-tuning
  coefficient, its oracle value, and its labeled-sample plug-in estimate. \\
  \(\widehat\theta^{(t)}\) & Power-tuned PPI mean estimator for task \(t\). \\
  \(Z_i,\overline Z_L,\overline Z_N\) & Residual \(Y_i-\gamma_t s_i\) and its
  labeled and population means. \\
  \(\Var_N,\Cov_N,\rho_N\) & Finite-population variance, covariance, and
  correlation, using denominator \(N\). \\
  \(\Var_L,\Cov_L\) & Variance and covariance computed on a task's labeled
  sample. \\
  \(S_Y^2,S_s^2,S_{Ys}\) & Finite-population sample variances of \(Y\) and
  \(s\), and their covariance, using denominator \(N-1\). \\
  \(V(s,\gamma_t),V^\star(s)\) & Design variance and oracle-minimized design
  variance. \\
  \(\mathcal H\) & Class of candidate recalibration maps. \\
  \(s^{(-t)}\) & Global recalibrator learned from labeled data in all
  tasks except \(t\). \\
  \(s_{\mathrm{local}}^{(t)}\) & Recalibrator learned from task-\(t\)
  labels without predicting an observation from its own outcome. \\
  \(\kappa_t,u_i^{(t)}\) & Weight on the local recalibrator and resulting
  adaptive proxy used by AREPPI. The earlier note reserves \(\gamma_t\) for
  power tuning, so \(\kappa_t\) replaces the source paper's adaptive-mixture
  symbol. \\
  \(\lambda,\gamma\) & Source-paper aliases for the power-tuning coefficient
  and AREPPI mixing weight; this note translates them to \(\gamma_t\) and
  \(\kappa_t\). \\
  \(\phi,\operatorname{id}\) & Generic transformation of \(\widehat Y\) and the
  identity transformation. \\
  \(\mathcal Z,I_z\) & Support of \(\widehat Y\) and the units whose proxy equals
  \(z\). \\
  \(m(z)\) & Finite-population conditional mean of \(Y\) at
  \(\widehat Y=z\). \\
  \(R_{Y\sim\widehat Y}^2\) & Fraction of finite-population outcome variance
  explained by the unrestricted conditional mean \(m(\widehat Y)\). \\
  \(a,b\) & Intercept and slope of an affine recalibration. \\
  \bottomrule
  \end{longtable}
  
  \noindent\textit{Notation conversion.}
  The earlier AI-project note denotes the power-tuning coefficient by
  \(\gamma_t\), whereas MT-PPI denotes it by \(\lambda\). This note keeps
  \(\gamma_t\) for consistency with the earlier note. It uses \(\kappa_t\) for
  AREPPI's local/global mixing weight, which MT-PPI denotes by \(\gamma\).
  
  \subsection*{Board 1: the many-task data structure}
  
  For each task \(t\in\mathcal T\), the paper fixes \(N\) units and targets
  \[
      \theta^{(t)}
      =
      \frac{1}{N}\sum_{i=1}^NY_i^{(t)}.
  \]
  The proxy \(\widehat Y_i^{(t)}\) is observed for every unit. Ground truth is
  observed only when \(i\in L^{(t)}\), where \(L^{(t)}\) is a simple random sample
  without replacement. (MT-PPI, pp.~3--4.)
  
  The target is one mean per task. The paper does not estimate one latent outcome
  per unit. Its vector of scientific estimands is
  \[
      \bigl(\theta^{(1)},\ldots,\theta^{(T)}\bigr).
  \]
  
  \subsection*{Board 2: power-tuned finite-population PPI}
  
  Suppress the task superscript. For a fixed recalibration map \(s\) and fixed
  coefficient \(\gamma_t\), define
  \[
      \widehat\theta^{(t)}
      =
      \overline Y_L+\gamma_t(\overline s_N-\overline s_L).
  \]
  This contains three familiar cases:
  \[
      \gamma_t=0
      \quad\Rightarrow\quad
      \widehat\theta^{(t)}=\overline Y_L,
  \]
  \[
      \gamma_t=1,\ s=\operatorname{id}
      \quad\Rightarrow\quad
      \text{ordinary PPI},
  \]
  and
  \[
      s=\operatorname{id},\ \gamma_t\ \text{estimated for precision}
      \quad\Rightarrow\quad
      \text{PPI++}.
  \]
  (MT-PPI, pp.~5--6.)
  
  Set \(Z_i=Y_i-\gamma_t s_i\). Subtract the target:
  \begin{align*}
      \widehat\theta^{(t)}-\theta^{(t)}
      &=
      (\overline Y_L-\gamma_t\overline s_L)
      -
      (\overline Y_N-\gamma_t\overline s_N)\\
      &=\overline Z_L-\overline Z_N.
  \end{align*}
  For fixed \(s\) and \(\gamma_t\), simple random sampling gives
  \[
      \E[\overline Z_L]=\overline Z_N,
  \]
  so \(\widehat\theta^{(t)}\) is exactly design-unbiased. Its variance is
  \[
      V(s,\gamma_t)
      =
      \frac{1}{n_t}
      \left(1-\frac{n_t}{N}\right)
      \left[
        S_Y^2-2\gamma_tS_{Ys}
        +\gamma_t^2S_s^2
      \right].
  \]
  Differentiate the bracket with respect to \(\gamma_t\):
  \[
      -2S_{Ys}+2\gamma_tS_s^2=0.
  \]
  Therefore
  \[
      \gamma^\star(s)
      =
      \frac{S_{Ys}}{S_s^2}
      =
      \frac{\Cov_N(Y,s)}{\Var_N(s)}.
  \]
  Substitution gives
  \[
      V^\star(s)
      =
      \frac{1}{n_t}\left(1-\frac{n_t}{N}\right)
      S_Y^2\{1-\rho_N^2(Y,s)\}.
  \]
  The best recalibrator is therefore the one most correlated with \(Y\), after
  allowing the scalar power coefficient to absorb linear scale. (MT-PPI,
  pp.~5--6 and 16.)
  
  \subsection*{Board 3: GREPPI}
  
  For target task \(t\), GREPPI proceeds in this order:
  \begin{enumerate}[leftmargin=2em]
      \item Pool labeled pairs from tasks other than \(t\).
      \item Fit \(s^{(-t)}\in\mathcal H\) on that pool.
      \item Apply \(s^{(-t)}\) to every proxy in task \(t\).
      \item Use task-\(t\) labels to estimate the task-\(t\) residual mean.
      \item Optionally estimate
      \[
          \widehat\gamma_t
          =
          \frac{\Cov_L(Y,s^{(-t)}(\widehat Y))}
               {\Var_L(s^{(-t)}(\widehat Y))}
      \]
      and form
      \[
          \widehat\theta^{(t)}
          =
          \overline Y_L^{(t)}
          +
          \widehat\gamma_t
          \left[
            \overline{s^{(-t)}}{}_N^{(t)}
            -
            \overline{s^{(-t)}}{}_L^{(t)}
          \right].
      \]
  \end{enumerate}
  The other tasks learn a better proxy. They do not supply the target task's
  bias correction. The target task's own labels still supply that correction.
  This is why cross-task learning can improve precision without requiring the
  tasks to have identical residual means. (MT-PPI, pp.~3 and 6.)
  
  If \(\gamma_t=1\), the leave-one-task-out recalibrator is fixed with respect to
  the target label draw, so the design-unbiasedness proof on Board 2 applies
  conditional on that recalibrator. If \(\widehat\gamma_t\) is estimated from
  the same target labels used for correction, exact finite-sample unbiasedness
  no longer follows. The authors describe the resulting inference as
  asymptotic and show small-sample undercoverage in their experiment. They
  therefore set the source paper's power coefficient to one in the human-data
  study. (MT-PPI, pp.~6 and 9--10.)
  
  \subsection*{Board 4: AREPPI}
  
  GREPPI can lose precision when other tasks have a different proxy--truth
  relationship. AREPPI combines a global and a local recalibrator:
  \[
      u_i^{(t)}
      =
      \kappa_t\,s_{\mathrm{local}}^{(t)}
          (\widehat Y_i^{(t)})
      +(1-\kappa_t)\,s^{(-t)}
          (\widehat Y_i^{(t)}).
  \]
  The local prediction for each labeled observation is out-of-fold: that
  observation's \(Y_i^{(t)}\) is not used to train the function that predicts it.
  The paper chooses \(\kappa_t\) to maximize squared correlation with target-task
  outcomes, then replaces \(s_i\) by \(u_i^{(t)}\) in the estimator on Board 2.
  (MT-PPI, p.~7.)
  
  Thus:
  \[
      \kappa_t=0
      \quad\Rightarrow\quad
      \text{use only the global cross-task relationship},
  \]
  while
  \[
      \kappa_t=1
      \quad\Rightarrow\quad
      \text{use only the local task relationship}.
  \]
  The adaptive mixture is a protection against task heterogeneity. It is not a
  license to omit target-task labels.
  
  \subsection*{Board 5: why a nonlinear relationship is necessary}
  
  First take an affine recalibration
  \[
      s(\widehat Y)=a\widehat Y+b,
      \qquad a\neq0.
  \]
  Then
  \[
      \Cov_N(Y,a\widehat Y+b)=a\Cov_N(Y,\widehat Y)
  \]
  and
  \[
      \Var_N(a\widehat Y+b)=a^2\Var_N(\widehat Y).
  \]
  Therefore
  \[
      \rho_N^2(Y,a\widehat Y+b)
      =
      \rho_N^2(Y,\widehat Y).
  \]
  Since \(V^\star(s)\) depends on \(s\) only through squared correlation,
  \[
      V^\star(a\widehat Y+b)=V^\star(\widehat Y).
  \]
  Power tuning already absorbs every shift and rescaling. Cross-task affine
  recalibration cannot improve oracle variance beyond power-tuned PPI.
  (MT-PPI, p.~8.)
  
  The stronger result follows in four steps. Define
  \[
      \mathcal Z=\{\widehat Y_i:i=1,\ldots,N\},
      \qquad
      I_z=\{i:\widehat Y_i=z\},
  \]
  and
  \[
      m(z)=\frac{1}{|I_z|}\sum_{i\in I_z}Y_i.
  \]
  For any transformation \(\phi\), observations with the same proxy value have
  the same \(\phi(\widehat Y_i)\). Hence
  \[
      \Cov_N\{Y-m(\widehat Y),\phi(\widehat Y)\}=0.
  \]
  It follows that
  \[
      \Cov_N\{Y,\phi(\widehat Y)\}
      =
      \Cov_N\{m(\widehat Y),\phi(\widehat Y)\}.
  \]
  Cauchy--Schwarz then gives
  \[
      \rho_N^2\{Y,\phi(\widehat Y)\}
      \leq
      \frac{\Var_N\{m(\widehat Y)\}}{\Var_N(Y)}
      =
      R_{Y\sim\widehat Y}^2.
  \]
  Choosing \(\phi=m\) attains the upper bound. The identity map attains it if
  and only if \(m\) is affine on the observed proxy support, apart from
  degenerate zero-variance cases. Therefore a strict improvement over
  power-tuned PPI is possible if and only if the proxy--truth conditional mean
  is nonlinear on that support. (MT-PPI, pp.~8 and 18--19.)
  
  \subsection*{Board 6: what MT-PPI adds and what it does not add}
  
  \begin{enumerate}[leftmargin=2em]
      \item It adds cross-task learning of a better, generally nonlinear proxy.
      \item It preserves task-specific correction: every reported task still
      uses labels from that task.
      \item It explains exactly when cross-task recalibration can improve on
      power-tuned PPI: only recoverable nonlinear structure matters.
      \item Its formal setup targets finite-population task means. The paper
      says extensions to general convex estimators are possible, but does not
      deliver the full extension. (MT-PPI, p.~3.)
      \item Estimating the power coefficient on very small target samples can
      cause optimistic variance estimates and undercoverage. Fixed
      \(\gamma_t=1\) or separate-sample tuning avoids the exact same-data
      argument. (MT-PPI, pp.~6 and 9--10.)
      \item It does not estimate a task with no target labels, because
      cross-task recalibration is not a substitute for within-task correction.
  \end{enumerate}
  
  \section{One design for the AI productivity project}
  
  \subsection*{Notation}
  
  \begin{longtable}{>{\raggedright\arraybackslash}p{0.22\linewidth}
                    >{\raggedright\arraybackslash}p{0.70\linewidth}}
  \toprule
  Symbol & Meaning \\
  \midrule
  \endfirsthead
  \toprule
  Symbol & Meaning \\
  \midrule
  \endhead
  \(j,J\) & Occupation--task cell and number of cells. \\
  \(i,N_j\) & Worker--task instance and number of target instances in cell \(j\). \\
  \(Y_{ij}\) & Gold-standard, quality-adjusted productivity gain for instance
  \(i\) in cell \(j\), after fixing one task basket. \\
  \(\widehat Y_{ij}\) & Cheap proxy gain observed for every target instance. \\
  \(\theta_j\) & Cell mean \(N_j^{-1}\sum_iY_{ij}\). \\
  \(L_j,n_j\) & Random labeled subset in cell \(j\) and its size. \\
  \(s^{(-j)}\) & Nonlinear recalibrator learned from other cells. \\
  \(\widehat\theta_j\) & Rectified estimate of the cell mean. \\
  \(p_j,\bar\theta\) & Aggregation weight and aggregate gain. \\
  \(p_j^0,p_j^A,\bar\theta^0,\bar\theta^A\) & The two basket-specific weights
  and their aggregate gains. \\
  \(S_j,g(S_j),\widehat v_j,R_j,\pi_j\) & Cell-level score, calibrated score
  prediction, unbiased experimental
  cell estimate, validation indicator, and validation inclusion probability in
  the alternative cell-sampling design. \\
  \(m\) & Number of validated cells in a simple random cell sample. \\
  \(k\) & Number of cells in a reported top-\(k\) set. \\
  \(\widehat{\bar\theta}_{\mathrm{PPI}}\) & Design-unbiased cell-level
  difference estimator of the aggregate. \\
  \(\widehat{\bar\theta}\) & Aggregate formed from the rectified cell estimates
  under the many-instances-per-cell design. \\
  \bottomrule
  \end{longtable}
  
  \subsection*{Board 1: first determine the data shape}
  
  There are two different designs.
  
  \paragraph{Design A: one score per cell.}
  The data contain \(S_j\) for every cell and an experimental
  \(\widehat v_j\) only for selected cells. This is the proposal's stated
  Gaussian-channel design. It supports:
  \begin{itemize}[leftmargin=2em]
      \item design-based estimation of aggregates from random validation;
      \item model-based estimation of unvalidated cell effects through empirical
      Bayes.
  \end{itemize}
  It does not fit MT-PPI's data structure, because there are not \(N_j\) proxy
  observations and some labels inside every task.
  
  \paragraph{Design B: many instances inside every cell.}
  After fixing one basket, the data contain \(\widehat Y_{ij}\) for all \(i\) in
  each cell and \(Y_{ij}\) for a random subset \(L_j\) in each cell. Then each
  cell is an
  MT-PPI task, and
  \[
      \theta_j=\frac{1}{N_j}\sum_{i=1}^{N_j}Y_{ij}
  \]
  is exactly a task mean. This design supports design-based estimation of every
  cell mean, provided every reported cell has target-cell labels.
  
  \subsection*{Board 2: the aggregate under Design A}
  
  Let validation include cell \(j\) with known probability \(\pi_j>0\), and let
  \(\widehat v_j\) be unbiased for \(\theta_j\). Treat \(g(S_j)\) as fixed with
  respect to cell \(j\)'s validation indicator and experimental noise. This is
  exact when \(g\) is specified before validation or trained on an independent
  calibration split. If the same validation observations train \(g\) and estimate
  the residual correction, sample splitting or design-compatible cross-fitting
  is required. Define
  \[
      \widehat{\bar\theta}_{\mathrm{PPI}}
      =
      \sum_{j=1}^Jp_jg(S_j)
      +
      \sum_{j:R_j=1}
      \frac{p_j}{\pi_j}
      \{\widehat v_j-g(S_j)\}.
  \]
  Take expectation over validation selection and experimental estimation:
  \begin{align*}
      \E[\widehat{\bar\theta}_{\mathrm{PPI}}]
      &=
      \sum_{j=1}^Jp_jg(S_j)
      +
      \sum_{j=1}^J
      \E[R_j]\frac{p_j}{\pi_j}
      \{\E[\widehat v_j]-g(S_j)\}\\
      &=
      \sum_{j=1}^Jp_jg(S_j)
      +
      \sum_{j=1}^J
      \pi_j\frac{p_j}{\pi_j}\{\theta_j-g(S_j)\}\\
      &=
      \sum_{j=1}^Jp_j\theta_j\\
      &=\bar\theta.
  \end{align*}
  This is a Horvitz--Thompson difference estimator, derived by applying the
  finite-population PPI logic to cells. It identifies the aggregate even though
  the individual unvalidated \(\theta_j\)'s remain unidentified without the
  proposal's channel model.
  
  For equal weights and a simple random sample of \(m\) out of \(J\) cells,
  \(\pi_j=m/J\), so the estimator becomes
  \[
      \widehat{\bar\theta}_{\mathrm{PPI}}
      =
      \frac{1}{J}\sum_{j=1}^Jg(S_j)
      +
      \frac{1}{m}\sum_{j:R_j=1}
      \{\widehat v_j-g(S_j)\}.
  \]
  This is the finite-population PPI mean estimator with cells as units.
  
  \subsection*{Board 3: the map under Design B}
  
  Fix one task basket. For each cell \(j\):
  \begin{enumerate}[leftmargin=2em]
      \item Use labeled pairs from other cells to learn
      \(s^{(-j)}\).
      \item Apply that map to every \(\widehat Y_{ij}\) in cell \(j\).
      \item Correct with the random target-cell labels:
      \[
          \widehat\theta_j
          =
          \frac{1}{N_j}\sum_{i=1}^{N_j}
          s^{(-j)}(\widehat Y_{ij})
          +
          \frac{1}{n_j}\sum_{i\in L_j}
          \left[
            Y_{ij}-s^{(-j)}(\widehat Y_{ij})
          \right].
      \]
      \item Conditional on the cross-cell recalibrator, simple random sampling
      within cell \(j\) makes the second term unbiased for the full-cell average
      residual. Therefore
      \[
          \E[\widehat\theta_j]=\theta_j.
      \]
      \item Aggregate after estimating the cells:
      \[
          \widehat{\bar\theta}
          =
          \sum_{j=1}^Jp_j\widehat\theta_j.
      \]
  \end{enumerate}
  This is GREPPI with occupation--task cells as tasks and worker--task instances
  as units. A nonlinear recalibrator can improve precision when the shared
  proxy--truth conditional mean is nonlinear. The within-cell correction, not
  cross-cell similarity, supplies design validity.
  
  \subsection*{Board 4: which method answers which project question}
  
  \begin{center}
  \begin{tabular}{>{\raggedright\arraybackslash}p{0.25\linewidth}
                  >{\raggedright\arraybackslash}p{0.30\linewidth}
                  >{\raggedright\arraybackslash}p{0.35\linewidth}}
  \toprule
  Question & Suitable method & Reason \\
  \midrule
  Weighted mean gain &
  Finite-population PPI difference estimator &
  Average validation residual corrects the average proxy error; the full map is
  not required. \\
  \addlinespace
  One mean for every cell, with labels inside every cell &
  MT-PPI/GREPPI, preferably with fixed or separately tuned power coefficient in
  very small samples &
  Other cells learn the nonlinear proxy transformation; each target cell
  supplies its own correction. \\
  \addlinespace
  One mean for every cell, but many cells have no labels &
  Proposal's empirical-Bayes channel, or another explicit cross-cell model &
  PPI and MT-PPI do not identify a label-free target cell. The map requires
  model-based extrapolation. \\
  \addlinespace
  Rankings and top-\(k\) cells &
  Posterior distribution for the full vector, followed by ranking uncertainty &
  A valid aggregate correction does not identify individual ranks. \\
  \addlinespace
  Cross-cell dispersion &
  An explicit second-moment procedure or full-vector model &
  The three documents do not establish that the mean estimator alone yields a
  valid dispersion estimator. \\
  \bottomrule
  \end{tabular}
  \end{center}
  
  \subsection*{Board 5: the clean project conclusion}
  
  \begin{enumerate}[leftmargin=2em]
      \item The proposal and the two PPI papers solve different statistical
      problems.
      \item The proposal targets a high-dimensional latent map and therefore
      needs a cross-cell model for unvalidated cells.
      \item PPI targets valid inference for a chosen functional and can estimate
      the aggregate without estimating the map.
      \item MT-PPI targets one mean in each of many tasks. It helps the project
      only if each occupation--task cell contains many proxy observations and
      at least some randomly sampled ground-truth observations.
      \item Cross-task recalibration should be nonlinear if power-tuned PPI is
      the benchmark; affine calibration adds no oracle precision.
      \item Randomization has two distinct roles:
      random cells support design-based aggregate inference, while random units
      within every cell support design-based cell means.
      \item Empirical Bayes and PPI are complements, not substitutes. PPI gives
      design-based correction for supported aggregates or cell means. Empirical
      Bayes supplies model-based shrinkage, extrapolation to unvalidated cells,
      and the joint distribution needed for ranks and top-\(k\) statements.
  \end{enumerate}
  
  \section*{References}
  
  \begin{enumerate}[leftmargin=2em]
      \item Angelopoulos, A. N., Bates, S., Fannjiang, C., Jordan, M. I., and
      Zrnic, T. (2023). \emph{Prediction-Powered Inference}.
      \href{https://arxiv.org/pdf/2301.09633}{arXiv:2301.09633v4}.
      \item Emmenegger, N., Stahler, E., and Podimata, C. (2026).
      \emph{Prediction-Powered Inference Across Many Tasks for AI Evaluation
      \& Social Science Research}.
      \href{https://arxiv.org/pdf/2605.29249}{arXiv:2605.29249v1}.
      \item \emph{The Distributional Effects of AI's Productivity Gains:
      A Data Fusion Approach}. Project proposal,
      \texttt{openai\_exchange\_proposal\_v4\_for\_ext.pdf}.
  \end{enumerate}
  
  \end{document}
